\scanpage{30}
\mspage{4/}

\deleted{\unsure{par}}en la normalisant \inserted{p. ex.} de telle façon que la longueur du cercle $|z|=1$ soit $\deleted{2\pi}1$, \margin{C'est la métrique déduite de \unsure{l'unité que|la métrique|l'unitaire} hermitienne canonique $z\bar z$ sur $\mathbb C$, via $z\mapsto\exp 2i\pi z$, de $\mathbb C$ sur $\mathbb C^*$.} on la définit sur toutes les $X\simeq\mathbb C^*$ du $=$ fait. D'ailleurs l'\unsure{anti-automorphisme|antiautomorphisme} $z\mapsto\bar z$ \deleted{\unsure{respecte}} conserve la mesure, donc si $X\simeq\mathbb C^*$, \uline{tt} autom. conforme respecte \deleted{\unsure{la}}la mesure canonique.

Il en résulte que dans le cas non orienté $X\simeq$ plan projectif réel moins $1$ pt, on peut descendre à partir de $X^{\mathrm{or}}$ et on trouver encore une métrique canonique plate, stable par tous les automorphismes conformes.

$\deleted{\unsure{3}}2''$) Cas $X\simeq\mathbb E^1_{\mathbb C}$. \unsure{``}Les métriques riemanniennes sur une \deleted{\unsure{vect}} droite vectorielle \unsure{1|$L$} complexe \inserted{$L$}, associées aux structures hermitiennes \inserted{déf. pos} sur $L$, (\deleted{\unsure{elles|celles}} elles sont toutes proportionnelles) ne sont \uline{pas} invariantes par autom \deleted{\unsure{de}}\inserted{hol} $L$, mais seulemt par le sous-groupe engendré par les translations, et les \deleted{\unsure{multipl.|autres}} homothéties par des $u\in\mathbb U=\{z\in\mathbb C\mid |z|=1\}$; elles sont quasi-invariantes par autom quelconques (lesquels sont \uline{affines}). Donc sur une $X\simeq\mathbb E^1_{\mathbb C}$, on peut définir, non \uline{une} métrique
% archive sheet number: 32
